\chapter{引论}
\label{chap:1}
刚体动力学是一门古老的学科，计算机的出现使其焕发了新生并被彻底改变。如今，动力学计算已经出现在计算机游戏、动画和虚拟现实软件、仿真器、运动控制系统以及各种工程设计与分析工具之中。在每一种场合，计算机所计算的都是物理系统的刚体近似在运动中所涉及的力、加速度等物理量。本书的主要目的是提供一组可在计算机上执行各种动力学计算的高效算法。书中对每个算法都作了详细描述，使读者既能理解它是如何工作的，又能理解它为什么高效；同时还解释了一些基本概念，如递推形式（recursive formulation）、分支诱导稀疏性（branch-induced sparsity）和关节体惯量（articulated-body inertia）。刚体动力学通常用三维向量来表述。然而，本书的主题是动力学算法，而这个主题更适合用六维向量来表述。因此，我们采用一种基于空间向量（spatial vector）的六维记法，这一记法将在第 2 章中介绍。与三维向量相比，空间记法大大减少了代数推导的篇幅，简化了描述和解释动力学算法的工作，也简化了在计算机上实现算法的过程。本章将对本书的主题作一简要介绍：简要谈谈动力学算法和空间向量，并说明本书的组织结构。1.1 动力学算法 刚体系统的动力学由其运动方程（equation of motion）来描述，运动方程给出了作用在系统上的力与其所产生的加速度之间的关系。动力学算法是一种计算与动力学相关的各个量的数值的程序。我们主要关注针对两类特定计算的算法：

正向动力学（forward dynamics）：计算给定刚体系统对给定外加力的加速度响应；逆动力学（inverse dynamics）：计算为使给定刚体系统产生给定的加速度响应而必须施加的力。正向动力学主要用于仿真。逆动力学则有多种用途，如：运动控制系统、轨迹规划、机械设计，以及作为正向动力学计算中的一个组成部分。在第 9 章中我们将讨论第三类计算，称为混合动力学（hybrid dynamics），其中一部分加速度和力变量为已知，任务是计算出其余的变量。刚体系统的运动方程可以写成如下标准形式：τ = H(q) ¨q + C(q, ˙q) . (1.1) 在这个方程中，q、˙q 和 ¨q 分别是位置、速度和加速度变量的向量，τ 是外加力的向量。H 是一个由惯性项组成的矩阵，写作 H(q) 是为了表明它是 q 的函数。C 是一个由力项组成的向量，用于计入科氏力和离心力、重力，以及除 τ 中各力之外作用于系统的所有其他的力。写作 C(q, ˙q) 是为了表明它同时依赖于 q 和 ˙q。H 和 C 合起来构成运动方程的系数，而 τ 和 ¨q 则是方程的变量。典型情况下，两个变量之一是给定的，另一个是待求的未知量。尽管习惯上写作 H(q) 和 C(q, ˙q)，但更准确的写法应当是 H(model, q) 和 C(model, q, ˙q)，以此表明 H 和 C 都依赖于它们所针对的具体刚体系统。符号 model 代表一组数据，它从组成部分的角度描述一个具体的刚体系统：刚体和关节的数量、它们相互连接的方式，以及与每个部件相关联的所有参数的取值（惯量参数、几何参数等等）。我们把这种描述称为系统模型（system model），以便将它与数学模型区分开来。区别在于：系统模型描述的是系统本身，而数学模型描述的是系统行为的某一方面。方程 1.1 就是一个数学模型。在计算机上，model 会是一个类型为“系统模型”的变量，其取值是一个数据结构，其中包含某个具体刚体系统的系统模型。让我们把正向动力学和逆动力学计算封装为一对函数 FD 和 ID。这两个函数满足 ¨q = FD(model, q, ˙q, τ) (1.2) 以及 τ = ID(model, q, ˙q, ¨q) . (1.3) 把这些方程与方程 1.1 相比较，显然 FD 和 ID 必须分别求得 H−1(τ −C) 和 H¨q +C。然而，

这两个方程的意义在于它们清晰地展示了每种计算的输入与输出。特别地，它们表明系统模型在这两种情况下都是输入。因此，人们期望实现 FD 和 ID 的算法能够适用于某一类刚体系统，并能利用系统模型中的数据求出该模型所描述的那个具体刚体系统的动力学。我们把按这种方式工作的算法称为基于模型的算法（model-based algorithm）。这种方法的一大优点是：只需编写、测试、形成文档并加以维护的一套计算机程序，就能计算一大类刚体系统中任意一个系统的动力学。我们所关注的两类主要系统称为运动树（kinematic tree）和闭环系统（closed-loop system）。粗略地说，运动树是不含运动环（kinematic loop）的任何刚体系统，而闭环系统是不是运动树的任何刚体系统。1 计算运动树的动力学比计算闭环系统的动力学要容易得多。基于模型的算法可以按照两个主要属性进行分类：它们执行什么计算，以及它们适用于哪一类系统。本书的主体由面向运动树和闭环系统的正向动力学算法与逆动力学算法组成。1.2 空间向量 三维空间中的刚体具有六个运动自由度；然而我们通常却用三维向量来表达它的动力学。于是，要写出刚体的运动方程，实际上必须写出两个向量方程：f = m aC 和 nC = I ˙ω + ω × Iω . (1.4) 第一个方程表示作用在刚体上的力与其质心线加速度之间的关系。第二个方程表示作用在刚体上、相对于其质心的力矩与其角加速度之间的关系。在空间向量记法中，我们使用把刚体运动的线性和角度两个方面结合在一起的六维向量。这样，线加速度与角加速度合并成一个空间加速度向量（spatial acceleration vector），力与力矩合并成一个空间力向量（spatial force vector），如此等等。采用空间记法后，刚体的运动方程可以写成 f = Ia + v ×∗Iv , (1.5) 其中 f 是作用在刚体上的空间力向量，v 和 a 分别是该刚体的空间速度和空间加速度，I 是该刚体的空间惯量张量（spatial inertia tensor）。符号 ×∗ 表示一种空间向量叉积运算。这个方程的一个明显特征1精确定义将在第 4 章中给出，其中包括对“运动环”的定义。

是它与方程 1.4 在某些符号的命名上存在冲突。解决这个问题的办法是在空间符号上方加上帽号（例如 ˆf）。空间记法所能带来的简化远不止运动方程这一处。例如，如果把两个刚体连接起来形成一个单一刚体，那么新刚体的空间惯量就由简单的公式 Inew = I1 + I2 给出，其中 I1 和 I2 是原先两个刚体的惯量。这个方程在三维向量方法中要由三个方程代替：一个用于计算新的质量，一个用于计算新的质心，还有一个用于计算新刚体对新质心的转动惯量。所有这些简化的总体效果是：与标准的三维向量记法相比，空间记法通常能把代数推导的篇幅至少压缩到原来的四分之一。一旦越过了代数这道障碍，分析者便可以更简洁地陈述问题，用更少的步骤求解，并得到更加紧凑的解。空间向量的另一个好处是它简化了编写计算机代码的过程：得到的代码更短，也更容易阅读、编写和调试，而其效率并不逊于使用三维向量的软件。例如，某些编程语言允许用一个形如 f = I*a + v.cross(I*v); 的语句来实现方程 1.5，其中每个变量的类型都已告知编译器。于是编译器知道 I 和 v 分别表示刚体惯量和空间运动向量（spatial motion vector），从而可以把表达式 I*v 编译成这样的代码：调用空间算术库中专为高效执行这种乘法而设计的例程。1.3 单位与记号 角度假定一律以弧度度量。除此之外，书中的方程与单位制的选取无关，只要求所选的单位制自身是一致的。在少数明确提到单位的地方，使用 SI 单位制。数学记号大体遵循 ISO 惯例：变量排成斜体；常数和函数排成正体；向量和矩阵排成粗斜体。向量用小写字母表示，2矩阵用大写字母表示。第 265 页附有一个简短的符号表，若干符号在索引中也设有条目。以下几个细节值得在此一提。符号 0 和 1 分别表示零矩阵和单位矩阵，0 也用来表示零向量。上标 −T 表示逆矩阵的转置，因此 A−T 即 (A−1)T。2方程 1.1 中的向量 C 是这条规则唯一的例外。

形如 a× 的表达式表示把 b 映射为 a × b 的算子。如果一个量被称为数组（array），那么它就是一个带编号的对象列表。若 λ 是数组，则它的第 i 个元素记作 λ(i)。如果一个符号带有前置上标（例如 Av），则该上标标识一个坐标系。坐标系也可以出现在下标中。空间向量有时以帽号标注，坐标向量则以下划线标注，这将在第 2 章中说明。1.4 读者指南 本书的内容可以分为三个部分：准备知识、主要算法和其他专题。准备性的章节是第 2、3、4 章；主要算法在第 5 至 8 章中讲述。具备足够背景知识的读者可能会发现自己可以直接跳读第 5 章。从第 9 章起，各章假定读者对先前的材料已有基本了解，除此之外它们均自成体系。第 2 章从第一性原理出发讲述空间向量代数；第 3 章讲解如何建立并分析刚体系统的运动方程。这两章是全书数学性最强的两章，而且相对于后续算法所需的最低限度而言，二者对各自论题的阐述都要深入得多。第 4 章描述系统模型的各个组成部分，后续章节所用到的许多量都在这里定义。第 5、6、7 章依次介绍运动树最著名的三种算法：用于逆动力学的递推牛顿–欧拉算法（recursive Newton-Euler algorithm），以及依此次序用于正向动力学的复合刚体算法（composite-rigid-body algorithm）与关节体算法（articulated-body algorithm）。第 5 章还解释了递推形式（recursive formulation）的概念，这正是这些算法之所以高效的根源。第 8 章随后给出几种计算闭环系统动力学的技术。这四章大体按复杂程度递增的顺序编排（即算法逐渐变得更复杂）。第 9、10、11 章继而处理各种补充专题。第 9 章考察另外几种算法，包括正向动力学与逆动力学的混合方法以及针对浮动基座（floating base）系统的算法；第 10 章讨论数值精度与计算效率问题；第 11 章考察受到接触与碰撞作用的刚体系统的动力学。读者会发现书中的例子散布各处且分布不均。有些例子用于阐释正文论及的概念；另一些例子则描述那些以例子的形式更容易讲清楚的想法。读者还会发现，每一个主要的算法既以一组方程的形式呈现，也以一个伪代码（pseudocode）程序的形式呈现，很多时候二者并排出现。伪代码力求一目了然，并且便于转写成真正的计算机代码。许多人——包括作者在内——都觉得在学习新算法时手头备有一些可运行的源代码会有帮助。本书不随书附带源代码，理由是这样的软件过时得太快。

读者应转而从网上寻找这些代码。在写作本书的时候，作者的个人网站3收录了此处所述绝大多数算法的源代码，不过读者或许也能在其他地方找到有用的软件。1.5 扩展阅读 近年来出版了多种关于刚体动力学的著作，只是其中不少如今已经绝版。适合作入门教材的有：Amirouche (2006)、Moon (1998)、Shabana (2001) 以及 Huston (1990)。前三本是配有大量例题与课后习题的正式教学用书，而 Huston 的那本则更具教程性质。所有这些书都面向本科生读者。更高阶的著作有 Stejskal 和 Val´aˇsek (1996)、Roberson 和 Schwertassek (1988) 以及 Wittenburg (1977)。想换个视角的读者不妨一试 Coutinho (2001)，它更多面向计算机游戏和虚拟现实领域。本书的许多内容源于机器人学界。其他关于机器人动力学的著作还有 Balafoutis 和 Patel (1991)、Lilly (1993)、Yamane (2004)，以及作者更早的作品 Featherstone (1987)——后者已被本书取代。在 Angeles (2003)、Khalil 和 Dombre (2002)、Siciliano 和 Khatib (2008) 以及其他多种机器人学书籍中，都可以找到有关动力学的节与章。对各种类型的六维向量怀有兴趣的读者，可能会觉得下列书籍和文章值得一读：Ball (1900)、Brand (1953)、von Mises (1924a,b)、Murray et al. (1994) 以及 Selig (1996)。此外，Featherstone (1987) 对空间向量的处理方式与本书所述略有不同。3该网址是（或者说曾经是）http://users.rsise.anu.edu.au/\textasciitilde{}roy/spatial/ 。如果这个网址打不开，就请尝试搜索“Roy Featherstone”。
